\documentclass[11pt,a4paper,sans]{moderncv}

% ----- CONFIGURACAO MODERNCV -----
\moderncvstyle{classic}
\moderncvcolor{blue}

% ----- PACOTES -----
\usepackage[utf8]{inputenc}
\usepackage[T1]{fontenc}
\usepackage[scale=0.84,top=1.2cm,bottom=1.2cm]{geometry}
\usepackage{lmodern}

% ----- DADOS PESSOAIS -----
\name{Luan Alfredo}{Lugarini}
\title{Estudante de Engenharia de Computação}
\address{Curitiba}{Paraná, Brasil}
% \phone[mobile]{+55 (41) 9XXXX-XXXX}  % descomente para a versao enviada por e-mail
\email{luanlugarini@gmail.com}
\social[github]{Luan1357}
\extrainfo{\href{http://lattes.cnpq.br/4564643932699945}{Currículo Lattes}}

\begin{document}

\makecvtitle

% ----- RESUMO -----
\small{Estudante do 4º período de Engenharia de Computação na PUCPR, em busca de estágio em engenharia. Tenho projetos acadêmicos em programação de sistemas em C (alocação dinâmica de memória, leitura de arquivos binários), Python e robótica (LEGO SPIKE e Arduino), incluindo um braço robótico montado e programado individualmente. Inglês avançado (C1).}
\vspace{4pt}

% ----- FORMACAO -----
\section{Formação}
\cventry{2025 -- 2029}{Bacharelado em Engenharia de Computação}{Pontifícia Universidade Católica do Paraná (PUCPR)}{Curitiba, PR}{4º período, conclusão prevista para 2029}{}
\cventry{2021 -- 2023}{Ensino Médio}{Colégio Imaculada Conceição}{}{}{}

% ----- PROJETOS -----
\section{Projetos Acadêmicos}
\cventry{2025}{Braço Robótico MyArm: Torre de Hanói}{Princípios de Robótica, PUCPR}{C++ / Arduino}{}{
\begin{itemize}
    \item Montei individualmente um braço robótico de 4 graus de liberdade (base, duas juntas e garra) controlado por servomotores e Arduino.
    \item Programei a sequência de movimentos que resolve fisicamente o jogo Torre de Hanói, com controle de velocidade dos servos para movimentos suaves.
    \item Código e vídeo: \href{https://github.com/Luan1357/myarm-torre-de-hanoi}{github.com/Luan1357/myarm-torre-de-hanoi}
\end{itemize}}
\cventry{2025}{Robô Autônomo para Labirinto}{Princípios de Robótica, PUCPR}{Python / LEGO SPIKE}{}{
\begin{itemize}
    \item Robô que percorre um labirinto físico sem intervenção, usando sensor de distância (paredes), sensor de cor (marcadores no chão) e giroscópio (curvas de 90° e 180° com correção de ângulo).
    \item Lógica organizada em máquina de estados (ida, volta e fim), com programação assíncrona.
    \item \href{https://github.com/Luan1357/robo-labirinto-lego}{github.com/Luan1357/robo-labirinto-lego}
\end{itemize}}
\cventry{2026}{Jogo da Vida de Conway}{Programação Imperativa, PUCPR}{C}{}{
\begin{itemize}
    \item Simulador do autômato celular com matriz alocada dinamicamente (\texttt{malloc}/\texttt{calloc}), leitura do estado inicial de arquivo e borda morta.
    \item Double buffering com troca de ponteiros entre gerações, evitando cópia de dados e atualização incorreta \emph{in-place}.
    \item \href{https://github.com/Luan1357/jogo-da-vida-conway}{github.com/Luan1357/jogo-da-vida-conway}
\end{itemize}}
\cventry{2026}{ppm2ascii: Conversor de Imagem para ASCII Art}{Programação Imperativa, PUCPR}{C}{}{
\begin{itemize}
    \item Leitura e validação completa de imagens PPM binárias (P6), com suporte a 8 e 16 bits por canal e verificação de overflow antes da alocação.
    \item Cálculo de luminância por pixel e mapeamento para uma rampa de 57 caracteres.
    \item \href{https://github.com/Luan1357/ppm2ascii}{github.com/Luan1357/ppm2ascii}
\end{itemize}}
\cventry{2026}{Jogo dos 15}{Programação Imperativa, PUCPR}{Python}{}{
\begin{itemize}
    \item Quebra-cabeça deslizante 4x4 no terminal, com embaralhamento que sempre gera um tabuleiro solucionável e resolução automática usando o histórico de jogadas como pilha.
    \item \href{https://github.com/Luan1357/jogo-dos-15}{github.com/Luan1357/jogo-dos-15}
\end{itemize}}

% ----- HABILIDADES -----
\section{Habilidades Técnicas}
\cvitem{Linguagens}{C, Python, C++ (Arduino)}
\cvitem{Hardware}{Arduino, servomotores, LEGO SPIKE (sensores de distância, cor e giroscópio)}
\cvitem{Ferramentas}{Git, GitHub, VS Code}

% ----- IDIOMAS -----
\section{Idiomas}
\cvitem{Português}{Nativo}
\cvitem{Inglês}{Avançado (C1), certificado pela Wise Up}

\end{document}
